\chapter{2011年山东卷（理）}

\section{单选题}

\begin{problem}
设集合 \(M = \{x \mid x^2 + x - 6 < 0\}\)，\(N = \{x \mid 1 \leqslant x \leqslant 3\}\). 则 \(M \cap N =\)（\quad）

\choices
    {\([1, 2)\)}
    {\([1, 2]\)}
    {\((2, 3]\)}
    {\([2, 3]\)}
\end{problem}

\begin{answer}
A
\end{answer}

\begin{solution}
由 \(x^{2}+x-6=(x+3)(x-2)\)，且二次项系数为正，得
\(x^{2}+x-6<0\) 的解集为 \(M=(-3,2)\). 又 \(N=[1,3]\)，所以
\(M\cap N=(-3,2)\cap[1,3]=[1,2)\)，故选 A.
\end{solution}

\begin{problem}
复数 \(z = \frac{2 - \i}{2 + \i}\)（\(\i\) 为虚数单位）在复平面内对应的点所在的象限为（\quad）

\choices
    {第一象限}
    {第二象限}
    {第三象限}
    {第四象限}
\end{problem}

\begin{answer}
D
\end{answer}

\begin{solution}
分子、分母同乘分母的共轭复数 \(2-\i\)，得
\(z=\dfrac{(2-\i)^2}{(2+\i)(2-\i)}=\dfrac{4-4\i+\i^2}{4-\i^2}=\dfrac{3-4\i}{5}=\dfrac35-\dfrac45\i\).
因此 \(z\) 的实部 \(\frac35>0\)，虚部 \(-\frac45<0\)，对应点在第四象限，故选 D.
\end{solution}

\begin{problem}
若点 \((a,9)\) 在函数 \(y = 3^{x}\) 的图象上，则 \(\tan \frac{a\pi}{6}\) 的值为（\quad）

\choices
    {0}
    {\( \frac{\sqrt{3}}{3} \)}
    {1}
    {\(\sqrt{3}\)}
\end{problem}

\begin{answer}
D
\end{answer}

\begin{solution}
点 \((a,9)\) 在 \(y=3^x\) 的图象上，所以 \(3^a=9=3^2\). 因为底数 \(3>1\)，指数函数为单调函数，故 \(a=2\). 于是
\(\tan\dfrac{a\pi}{6}=\tan\dfrac{\pi}{3}=\sqrt3\)，故选 D.
\end{solution}

\begin{problem}
不等式 \( |x - 5| + |x + 3| \geqslant 10 \) 的解集是（\quad）

\choices
    {\([-5,7]\)}
    {\([-4,6]\)}
    {\((- \infty, -5] \cup [7, +\infty)\)}
    {\((- \infty, -4] \cup [6, +\infty)\)}
\end{problem}

\begin{answer}
D
\end{answer}

\begin{solution}
按 \(-3\)、\(5\) 分段讨论. 当 \(x<-3\) 时，
\(|x-5|+|x+3|=(5-x)+(-x-3)=2-2x\)，故 \(2-2x\geqslant10\)，得 \(x\leqslant-4\).
当 \(-3\leqslant x\leqslant5\) 时，
\(|x-5|+|x+3|=(5-x)+(x+3)=8<10\)，没有满足条件的 \(x\).
当 \(x>5\) 时，
\(|x-5|+|x+3|=(x-5)+(x+3)=2x-2\)，故 \(2x-2\geqslant10\)，得 \(x\geqslant6\).
所以解集为 \((-\infty,-4]\cup[6,+\infty)\)，故选 D.
\end{solution}

\begin{problem}
对于函数 \(y = f(x)\)，\(x \in \R\). “\( y = |f(x)| \) 的图象关于 \( y \) 轴对称”是“\( y = f(x) \) 是奇函数”的（\quad）

\choices
    {充分而不必要条件}
    {必要而不充分条件}
    {充要条件}
    {既不充分也不必要条件}
\end{problem}

\begin{answer}
B
\end{answer}

\begin{solution}
若 \(f\) 是奇函数，则对任意 \(x\in\R\) 有 \(f(-x)=-f(x)\)，从而
\(|f(-x)|=|-f(x)|=|f(x)|\). 这正说明 \(y=|f(x)|\) 的图象关于 \(y\) 轴对称，因此题目前一个条件是必要条件.
但反过来不成立. 例如取 \(f(x)=1\)，则 \(|f(x)|=1\) 的图象关于 \(y\) 轴对称，而 \(f(-x)=1\ne-1=-f(x)\)，所以 \(f\) 不是奇函数.
故前一个条件是后一个条件的必要而不充分条件，选 B.
\end{solution}

\begin{problem}
若函数 \( f(x) = \sin \omega x \)（\(\omega > 0\)）在区间 \(\left[0, \frac{\pi}{3}\right]\) 上单调递增，在区间 \(\left[\frac{\pi}{3}, \frac{\pi}{2}\right]\) 上单调递减，则 \(\omega =\)（\quad）

\choices
    {3}
    {2}
    {\(\frac{3}{2}\)}
    {\(\frac{2}{3}\)}
\end{problem}

\begin{answer}
C
\end{answer}

\begin{solution}
\(f'(x)=\omega\cos(\omega x)\)，且 \(\omega>0\). 逐一检验选项：
当 \(\omega=3\) 时，\(f'(\frac{\pi}{3})=3\cos\pi<0\)，不可能在第一段上单调递增；
当 \(\omega=2\) 时，\(f'(\frac{\pi}{3})=2\cos\frac{2\pi}{3}<0\)，同样不满足第一段要求；
当 \(\omega=\frac23\) 时，\(f'(\frac{\pi}{2})=\frac23\cos\frac{\pi}{3}>0\)，不可能在第二段上单调递减.
当 \(\omega=\frac32\) 时，\(\omega x\) 在 \([0,\frac{\pi}{3}]\) 上取值于 \([0,\frac\pi2]\)，故 \(f'(x)\geqslant0\)；在 \([\frac\pi3,\frac\pi2]\) 上取值于 \([\frac\pi2,\frac{3\pi}4]\)，故 \(f'(x)\leqslant0\). 条件满足，因此 \(\omega=\frac32\)，选 C.
\end{solution}

\begin{problem}
某产品的广告费用 \(x\) 与销售额 \(y\) 的统计数据如下表：

\begin{center}
\begin{tabular}{c|cccc}
广告费用 \(x\) （万元） & 4 & 2 & 3 & 5 \\
\hline
销售额 \(y\) （万元） & 49 & 26 & 39 & 54
\end{tabular}
\end{center}

根据上表可得回归方程 \( \hat{y} = \hat{b}x + \hat{a} \) 中的 \(\hat{b}\) 为 \(9.4\)，据此模型预报广告费用为 \(6\text{ 万元}\) 时销售额为（\quad）

\choices
    {\(63.6\text{ 万元}\)}
    {\(65.5\text{ 万元}\)}
    {\(67.7\text{ 万元}\)}
    {\(72.0\text{ 万元}\)}
\end{problem}

\begin{answer}
B
\end{answer}

\begin{solution}
样本均值为 \(\bar{x}=\dfrac{4+2+3+5}{4}=\dfrac72\)，
\(\bar{y}=\dfrac{49+26+39+54}{4}=42\). 最小二乘回归直线经过样本中心 \((\bar{x},\bar{y})\)，所以
\(\hat a=\bar y-\hat b\bar x=42-9.4\times\dfrac72=9.1\).
当广告费用为 \(6\text{ 万元}\) 时，
\(\hat y=9.4\times6+9.1=65.5\text{ 万元}\)，故选 B.
\end{solution}

\begin{problem}
已知双曲线 \(\frac{x^2}{a^2} - \frac{y^2}{b^2} = 1 \) （\(a > 0\)，\(b > 0\)）的两条渐近线均和圆 \(C: x^2 + y^2 - 6x + 5 = 0\) 相切，且双曲线右焦点为圆 \(C\) 的圆心，则该双曲线的方程为（\quad）

\choices
    {\(\frac{x^2}{5} -\frac{y^2}{4} = 1\)}
    {\(\frac{x^2}{4} -\frac{y^2}{5} = 1\)}
    {\(\frac{x^2}{3} -\frac{y^2}{6} = 1\)}
    {\(\frac{x^2}{6} -\frac{y^2}{3} = 1\)}
\end{problem}

\begin{answer}
A
\end{answer}

\begin{solution}
将圆方程配方得 \((x-3)^2+y^2=4\)，所以圆心为 \(O_C=(3,0)\)，半径为 \(2\). 双曲线右焦点为 \(O_C\)，故焦距参数 \(c=3\)，从而
\(a^2+b^2=c^2=9\).
双曲线的渐近线为 \(y=\pm\dfrac ba x\). 以 \(y=\dfrac ba x\) 为例，圆心到该直线的距离为
\(\dfrac{3b/a}{\sqrt{1+b^2/a^2}}=\dfrac{3b}{\sqrt{a^2+b^2}}=b\).
渐近线与圆相切，所以该距离等于半径，得 \(b=2\)，进而 \(a^2=9-4=5\). 因此双曲线方程为 \(\dfrac{x^2}{5}-\dfrac{y^2}{4}=1\)，故选 A.
\end{solution}

\begin{problem}
函数 \( y = \frac{x}{2} - 2\sin x \) 的图象大致是（\quad）

\choices
    {\choicebitmap[width=0.15\paperwidth]{img/2011/shandong_science/q9_choice_a.png}}
    {\choicebitmap[width=0.15\paperwidth]{img/2011/shandong_science/q9_choice_b.png}}
    {\choicebitmap[width=0.15\paperwidth]{img/2011/shandong_science/q9_choice_c.png}}
    {\choicebitmap[width=0.15\paperwidth]{img/2011/shandong_science/q9_choice_d.png}}
\end{problem}

\begin{answer}
C
\end{answer}

\begin{solution}
记 \(g(x)=\dfrac{x}{2}-2\sin x\). 则
\(g(0)=0\)，且 \(g'(x)=\dfrac12-2\cos x\)，所以 \(g'(0)=g'(2\pi)=-\dfrac32<0\). 因此图象在原点处从上向下穿过 \(x\) 轴，并且在 \(x=2\pi\) 处于下降状态.
另外，\(g(2\pi)=\pi>0\)，而 \(g(x+2\pi)=g(x)+\pi\)，说明每相隔 \(2\pi\) 的波动形状相同但整体上移；同时 \(g(-x)=-g(x)\)，图象关于原点中心对称. 只有选项 C 同时符合这些性质，故选 C.
\end{solution}

\begin{problem}
已知 \( f(x) \) 是 \( \R \) 上最小正周期为 2 的周期函数，且当 \( 0 \leqslant x < 2 \) 时，\( f(x) = x^3 - x \)，则函数 \( y = f(x) \) 的图象在区间 \([0, 6]\) 上与 \( x \) 轴的交点的个数为（\quad）

\choices
    {6}
    {7}
    {8}
    {9}
\end{problem}

\begin{answer}
B
\end{answer}

\begin{solution}
在 \(0\leqslant x<2\) 时，
\(f(x)=x^3-x=x(x-1)(x+1)\). 在该区间内零点为 \(x=0\)，\(1\).
由周期性 \(f(x+2)=f(x)\)，所以在每个区间 \([2k,2k+2)\) 内零点为 \(2k\) 和 \(2k+1\). 对 \(k=0\)，\(1\)，\(2\)，得到 \(0\)，\(1\)，\(2\)，\(3\)，\(4\)，\(5\)；闭区间还包含端点 \(6\)，且 \(f(6)=f(0)=0\). 因此交点横坐标为 \(0\)，\(1\)，\(2\)，\(3\)，\(4\)，\(5\)，\(6\)，共有 \(7\) 个，故选 B.
\end{solution}

\begin{problem}
如图是长和宽分别相等的两个矩形. 给定下列三个命题：

\begin{circlelist}
\item 存在三棱柱，其正（主）视图，俯视图如图；
\item 存在四棱柱，其正（主）视图，俯视图如图；
\item 存在圆柱，其正（主）视图，俯视图如图.
\end{circlelist}

其中真命题的个数是（\quad）
\bitmapfigure[max width=0.44\linewidth]{img/2011/shandong_science/q11_fig1.png}

\choices
    {3}
    {2}
    {1}
    {0}
\end{problem}

\begin{answer}
A
\end{answer}

\begin{solution}
设图中两个矩形的长、宽分别为 \(l\)、\(w\). 对三种几何体分别构造如下.
取一个长为 \(l\)、宽和高均为 \(w\) 的长方体，则从两个相互垂直的方向观察，正视图和俯视图都为 \(l\times w\) 的矩形，所以第二个命题为真.
取轴长为 \(l\) 的直三棱柱，使其垂直于轴的三角形截面为直角边长均为 \(w\) 的等腰直角三角形，并令两条直角边分别沿正视和俯视方向. 该三角形在这两个方向上的投影都是长为 \(w\) 的线段，因此两个视图都为 \(l\times w\) 的矩形，第一个命题为真.
取轴长为 \(l\)、底面直径为 \(w\) 的直圆柱，两个视图分别是轴长为 \(l\)、直径为 \(w\) 的矩形，第三个命题也为真.
所以三个命题都是真命题，真命题有 \(3\) 个，故选 A.
\end{solution}

\begin{problem}
设 \(A_{1}\)，\(A_{2}\)，\(A_{3}\)，\(A_{4}\) 是平面直角坐标系中两两不同的四点. 若 \(\overrightarrow{A_{1}A_{3}} = \lambda \overrightarrow{A_{1}A_{2}}\) （\(\lambda \in \R\)），\(\overrightarrow{A_{1}A_{4}} = \mu \overrightarrow{A_{1}A_{2}}\) （\(\mu \in \R\)），且 \(\frac{1}{\lambda} + \frac{1}{\mu} = 2\)，则称 \(A_{3}\)，\(A_{4}\) 调和分割 \(A_{1}\)，\(A_{2}\). 已知平面上的点 \(C\)，\(D\) 调和分割点 \(A\)，\(B\)，则下面说法正确的是（\quad）

\choices
    {\(C\) 可能是线段 \(AB\) 的中点}
    {\(D\) 可能是线段 \(AB\) 的中点}
    {\(C\)，\(D\) 可能同时在线段 \(AB\) 上}
    {\(C\)，\(D\) 不可能同时在线段 \(AB\) 的延长线上}
\end{problem}

\begin{answer}
D
\end{answer}

\begin{solution}
设 \(\overrightarrow{AC}=\lambda\overrightarrow{AB}\)、\(\overrightarrow{AD}=\mu\overrightarrow{AB}\). 由于四点两两不同，\(\lambda\)、\(\mu\) 均不等于 \(0\) 或 \(1\). 由
\(\dfrac1\lambda+\dfrac1\mu=2\) 得
\(\mu=\dfrac{\lambda}{2\lambda-1}\)，其中 \(\lambda\ne\frac12\).
若 \(C\) 是线段 \(AB\) 的中点，则 \(\lambda=\frac12\)，代入原关系会得到 \(\frac1\mu=0\)，不可能；同理 \(D\) 也不可能是中点，所以前两项均错误.
若 \(0<\lambda<\frac12\)，则 \(\mu<0\)；若 \(\frac12<\lambda<1\)，则 \(\mu=\frac{\lambda}{2\lambda-1}>1\). 因而当 \(C\) 在线段 \(AB\) 上时，\(D\) 在线段外，不能同时在线段上，第三项错误.
若 \(\lambda<0\)，则 \(0<\mu<\frac12\)；若 \(\lambda>1\)，则 \(\frac12<\mu<1\). 因而当 \(C\) 在线段 \(AB\) 的延长线上时，\(D\) 反而在线段 \(AB\) 内，二者不可能同时在延长线上. 故只有 D 正确，选 D.
\end{solution}

\section{填空题}

\begin{problem}
执行如图所示的程序框图，输入 \(l = 2\)，\(m = 3\)，\(n = 5\)，则输出的 \(y\) 的值是\fillinblank{}.

\bitmapfigure[width=6.5cm]{img/2011/shandong_science/q13_fig1.png}
\end{problem}

\begin{answer}
68
\end{answer}

\begin{solution}
输入 \(l=2\)、\(m=3\)、\(n=5\) 后，先计算
\(l^2+m^2+n^2=2^2+3^2+5^2=38\ne0\)，所以程序进入“否”分支，令
\(y=70l+21m+15n=70\times2+21\times3+15\times5=278\).
因为 \(278>105\)，循环第一次执行 \(y\leftarrow y-105\)，得到 \(173\)；此时仍有 \(173>105\)，第二次执行后得到 \(68\). 由于 \(68\leqslant105\)，循环结束，输出 \(y=68\).
\end{solution}

\begin{problem}
若 \(\left(x - \frac{\sqrt{a}}{x^{2}}\right)^6\) 展开式的常数项为 60，则常数 \(a\) 的值为\fillinblank{}.
\end{problem}

\begin{answer}
4
\end{answer}

\begin{solution}
展开式中取第二项 \(k\) 次的项为
\(\mathrm C_6^k x^{6-k}\left(-\dfrac{\sqrt a}{x^2}\right)^k
=\mathrm C_6^k(-\sqrt a)^k x^{6-3k}\).
要成为常数项，必须有 \(6-3k=0\)，所以 \(k=2\). 此时常数项为
\(\mathrm C_6^2(-\sqrt a)^2=\mathrm C_6^2a=15a\).
由题意 \(15a=60\)，得 \(a=4\).
\end{solution}

\begin{problem}
设函数 \(f(x) = \frac{x}{x + 2}\)（\(x > 0\)），观察：\(f_{1}(x) = f(x) = \frac{x}{x + 2}\)，\(f_{2}(x) = f\!\left(f_{1}(x)\right) = \frac{x}{3x + 4}\)，\(f_{3}(x) = f\!\left(f_{2}(x)\right) = \frac{x}{7x + 8}\)，\(f_{4}(x) = f\!\left(f_{3}(x)\right) = \frac{x}{15x + 16}\). 根据以上事实，由归纳推理可得：当 \(n \in \N^{*}\) 且 \(n \geqslant 2\) 时，\(f_n(x) = f(f_{n-1}(x)) =\)\fillinblank{}.
\end{problem}

\begin{answer}
\(\displaystyle \frac{x}{(2^{n}-1)x+2^{n}}\)
\end{answer}

\begin{solution}
猜想
\(f_n(x)=\dfrac{x}{(2^n-1)x+2^n}\). 当 \(n=2\) 时，题中给出的
\(f_2(x)=\dfrac{x}{3x+4}=\dfrac{x}{(2^2-1)x+2^2}\)，猜想成立.
假设当 \(n=k-1\) 时结论成立，即
\(f_{k-1}(x)=\dfrac{x}{(2^{k-1}-1)x+2^{k-1}}\). 因为 \(x>0\)，分母非零，故
\[
f_k(x)=f\left(f_{k-1}(x)\right)
=\frac{\dfrac{x}{(2^{k-1}-1)x+2^{k-1}}}
{\dfrac{x}{(2^{k-1}-1)x+2^{k-1}}+2}
=\frac{x}{x+2\bigl((2^{k-1}-1)x+2^{k-1}\bigr)}
=\frac{x}{(2^k-1)x+2^k}
\]
所以由数学归纳法，对所有 \(n\geqslant2\) 都有
\(f_n(x)=\dfrac{x}{(2^n-1)x+2^n}\).
\end{solution}

\begin{problem}
已知函数 \( f(x) = \log_a x + x - b \)（\(a > 0\)，且 \(a \neq 1\)）. 当 \( 2 < a < 3 < b < 4 \) 时，函数 \( f(x) \) 的零点 \(x_0 \in (n, n+1)\)，\(n \in \N^{*}\)，则 \( n = \)\fillinblank{}.
\end{problem}

\begin{answer}
2
\end{answer}

\begin{solution}
由 \(2<a<3\) 且 \(a>1\)，有 \(\log_a2<1\)，再由 \(b>3\) 得
\(f(2)=\log_a2+2-b<1+2-3=0\).
同理，\(a<3\) 且 \(a>1\) 时 \(\log_a3>1\)，再由 \(b<4\) 得
\(f(3)=\log_a3+3-b>1+3-4=0\).
函数 \(f\) 在 \((0,+\infty)\) 上连续，且
\(f'(x)=\dfrac1{x\ln a}+1>0\)，所以它严格递增. 由介值定理，零点存在于 \((2,3)\)；由严格递增性，零点唯一. 因而题设中的 \(x_0\in(n,n+1)\) 满足 \((n,n+1)=(2,3)\)，故 \(n=2\).
\end{solution}

\section{解答题}

\begin{problem}
在 \(\triangle ABC\) 中，内角 \(A\)，\(B\)，\(C\) 的对边分别为 \(a\)，\(b\)，\(c\)，已知 \(\frac{\cos A - 2 \cos C}{\cos B} = \frac{2c - a}{b}\).
\begin{enumerate}
\item 求 \(\frac{\sin C}{\sin A}\) 的值；
\item 若 \(\cos B = \frac{1}{4}\)，\(b = 2\)，求 \(\triangle ABC\) 的面积 \(S\).
\end{enumerate}
\end{problem}

\begin{solution}
\begin{enumerate}
\item 由正弦定理，原条件等价于
\(\dfrac{\cos A-2\cos C}{\cos B}=\dfrac{2\sin C-\sin A}{\sin B}\). 由于原式有意义，\(\cos B\ne0\)，且三角形内角的正弦均为正，交叉相乘并整理得
\(\sin B\cos A+\cos B\sin A=2(\sin B\cos C+\cos B\sin C)\).
于是 \(\sin(A+B)=2\sin(B+C)\). 由 \(A+B+C=\pi\)，有
\(\sin(A+B)=\sin C\)、\(\sin(B+C)=\sin A\)，所以 \(\sin C=2\sin A\)，从而
\(\dfrac{\sin C}{\sin A}=2\).
\item 由第一问和正弦定理，\(\dfrac ca=\dfrac{\sin C}{\sin A}=2\)，所以 \(c=2a\). 在角 \(B\) 所对的边上使用余弦定理：
\(b^2=a^2+c^2-2ac\cos B=a^2+4a^2-2a\cdot2a\cdot\dfrac14=4a^2\).
因 \(a>0\)、\(b=2\)，得 \(a=1\)，进而 \(c=2\). 又 \(B\in(0,\pi)\)，所以
\(\sin B=\sqrt{1-\cos^2B}=\sqrt{1-\dfrac1{16}}=\dfrac{\sqrt{15}}4\).
故三角形面积为
\(S=\dfrac12ac\sin B=\dfrac12\cdot1\cdot2\cdot\dfrac{\sqrt{15}}4=\dfrac{\sqrt{15}}4\).
\end{enumerate}
\end{solution}

\begin{problem}
红队队员甲，乙，丙与蓝队队员 \(A\)，\(B\)，\(C\) 进行围棋比赛，甲对 \(A\)，乙对 \(B\)，丙对 \(C\) 各一盘. 已知甲胜 \(A\)，乙胜 \(B\)，丙胜 \(C\) 的概率分别为 \(0.6\)，\(0.5\)，\(0.5\). 假设各盘比赛结果相互独立.
\begin{enumerate}
\item 求红队至少两名队员获胜的概率；
\item 用 \(\xi\) 表示红队队员获胜的总数，求 \(\xi\) 的分布列和数学期望 \(E(\xi)\).
\end{enumerate}
\end{problem}

\begin{solution}
\begin{enumerate}
\item 红队恰有两名队员获胜包括三种互斥情形：甲、丙胜而乙负，甲、乙胜而丙负，乙、丙胜而甲负. 由独立性，概率为
\(0.6\times0.5\times0.5+0.6\times0.5\times0.5+0.4\times0.5\times0.5=0.40\).
三名队员全部获胜的概率为 \(0.6\times0.5\times0.5=0.15\). 因此红队至少两名队员获胜的概率为 \(0.40+0.15=0.55\).
\item 设三场比赛的胜负分别由甲、乙、丙表示，则
\(\mathrm P(\xi=0)=0.4\times0.5\times0.5=0.10\)，
\(\mathrm P(\xi=1)=0.6\times0.5\times0.5+0.4\times0.5\times0.5+0.4\times0.5\times0.5=0.35\)，
\(\mathrm P(\xi=2)=0.40\)，
\(\mathrm P(\xi=3)=0.6\times0.5\times0.5=0.15\).
所以 \(\xi\) 的分布列为
\begin{center}
\begin{tabular}{c|cccc}
\(\xi\) & 0 & 1 & 2 & 3 \\
\hline
\(P\) & \(0.10\) & \(0.35\) & \(0.40\) & \(0.15\)
\end{tabular}
\end{center}
因此 \(E(\xi)=0\times0.10+1\times0.35+2\times0.40+3\times0.15=1.60\).
\end{enumerate}
\end{solution}

\begin{problem}
在如图所示的几何体中，四边形 \(ABCD\) 为平行四边形，\(\angle ACB = 90^{\circ}\)，\(EA \perp\) 平面 \(ABCD\)，\(EF \parallel AB\)，\(FG \parallel BC\)，\(EG \parallel AC\)，\(AB = 2EF\).
\begin{enumerate}
\item 若 \(M\) 是线段 \(AD\) 的中点，求证：\(GM \parallel\) 平面 \(ABFE\)；
\item 若 \(AC = BC = 2AE\)，求二面角 \(A - BF - C\) 的大小.
\end{enumerate}

\bitmapfigure[max width=0.44\linewidth]{img/2011/shandong_science/q19_fig1.png}
\end{problem}

\begin{solution}
\begin{enumerate}
\item 设 \(\overrightarrow{AB}=\bs b\)、\(\overrightarrow{BC}=\bs d\)、\(\overrightarrow{AE}=\bs e\). 由图示的对应方向以及 \(AB=2EF\)，三角形 \(EFG\) 与三角形 \(ABC\) 相似，比例为 \(\frac12\)，故
\(\overrightarrow{EF}=\frac12\bs b\)，\(\overrightarrow{FG}=\frac12\bs d\).
于是 \(F=E+\frac12\bs b\)，\(G=E+\frac12(\bs b+\bs d)\). 又因为 \(ABCD\) 是平行四边形，\(\overrightarrow{AD}=\overrightarrow{BC}=\bs d\)，所以中点 \(M\) 满足 \(M=A+\frac12\bs d\). 因而
\(\overrightarrow{GM}=M-G=-\bs e-\frac12\bs b\).
向量 \(\overrightarrow{GM}\) 是 \(\bs b\)、\(\bs e\) 的线性组合，而平面 \(ABFE\) 的方向向量正由 \(\bs b\)、\(\bs e\) 张成，所以 \(GM\) 的方向平行于平面 \(ABFE\). 又 \(M\) 在底面 \(ABCD\) 内但不在直线 \(AB\) 上，且平面 \(ABFE\) 与底面的交线为 \(AB\)，故 \(M\notin ABFE\). 因此 \(GM\parallel\) 平面 \(ABFE\).
\item 令 \(AE=h>0\). 在底面内取 \(AC\) 为一个坐标轴，并建立空间直角坐标系，使
\(A=(0,0,0)\)、\(C=(2h,0,0)\)、\(B=(2h,2h,0)\)、\(E=(0,0,h)\). 此时
\(\overrightarrow{CA}=(-2h,0,0)\)、\(\overrightarrow{CB}=(0,2h,0)\)，确有 \(\angle ACB=90^\circ\)，且 \(AC=BC=2h=2AE\).
由 \(\overrightarrow{EF}=\frac12\overrightarrow{AB}\)，得 \(F=(h,h,h)\). 所以
\(\overrightarrow{BA}=(-2h,-2h,0)\)，\(\overrightarrow{BF}=(-h,-h,h)\)，\(\overrightarrow{BC}=(0,-2h,0)\).
平面 \(ABF\)、\(CBF\) 的法向量可分别取
\(\bs n_1=\overrightarrow{BA}\times\overrightarrow{BF}=(-2h^2,2h^2,0)\)，
\(\bs n_2=\overrightarrow{BC}\times\overrightarrow{BF}=(-2h^2,0,-2h^2)\).
于是
\(\cos\theta=\dfrac{|\bs n_1\cdot\bs n_2|}{|\bs n_1|\,|\bs n_2|}=\dfrac{4h^4}{(2\sqrt2h^2)(2\sqrt2h^2)}=\frac12\).
二面角 \(A-BF-C\) 取两平面的较小夹角，故 \(\theta=60^\circ\).
\end{enumerate}
\end{solution}

\begin{problem}
等比数列 \(\{a_{n}\}\) 中，\(a_{1}\)，\(a_{2}\)，\(a_{3}\) 分别是下表第一，二，三行中的某一个数，且 \(a_{1}\)，\(a_{2}\)，\(a_{3}\) 中的任何两个数不在下表的同一列.

\begin{center}
\begin{tabular}{c|ccc}
 & 第一列 & 第二列 & 第三列 \\
\hline
第一行 & 3 & 2 & 10 \\
第二行 & 6 & 4 & 14 \\
第三行 & 9 & 8 & 18
\end{tabular}
\end{center}

\begin{enumerate}
\item 求数列 \(\{a_{n}\}\) 的通项公式；
\item 若数列 \(\{b_{n}\}\) 满足： \(b_{n} = a_{n} + (-1)^{n} \ln a_{n}\)，求数列 \(\{b_{n}\}\) 的前 \(n\) 项和 \(S_{n}\).
\end{enumerate}
\end{problem}

\begin{solution}
\begin{enumerate}
\item 设公比为 \(q\). 因为 \(a_3=a_2q\)、\(a_2=a_1q\)，所以
\(a_3=\dfrac{a_2^2}{a_1}\). 逐一检验第一行和第二行的组合：
\(a_1=3\) 时，\(a_2=6\)，\(4\)，\(14\) 分别要求 \(a_3=12\)，\(\frac{16}{3}\)，\(\frac{196}{3}\)，均不在第三行；
\(a_1=2\) 时，\(a_2=6\) 要求 \(a_3=18\)，三项所在列依次为第二、第一、第三列，符合条件；\(a_2=4\) 时 \(a_3=8\)，与 \(a_1\) 同列；\(a_2=14\) 时 \(a_3=98\)，不在第三行；
\(a_1=10\) 时，\(a_2=6\)，\(4\)，\(14\) 分别要求 \(a_3=\frac{18}{5}\)，\(\frac85\)，\(\frac{98}{5}\)，均不在第三行. 因此唯一可能的三项为 \(a_1=2\)、\(a_2=6\)、\(a_3=18\). 公比
\(q=\dfrac{a_2}{a_1}=3\)，故通项公式为 \(a_n=2\cdot3^{n-1}\).
\item 因为 \(a_n=2\cdot3^{n-1}>0\)，所以
\(b_n=2\cdot3^{n-1}+(-1)^n\bigl(\ln2+(n-1)\ln3\bigr)\). 于是
\(S_n=\sum_{k=1}^n2\cdot3^{k-1}+\ln2\sum_{k=1}^n(-1)^k+\ln3\sum_{k=1}^n(-1)^k(k-1)\).
其中 \(\sum_{k=1}^n2\cdot3^{k-1}=3^n-1\). 若 \(n\) 为偶数，则
\(\sum_{k=1}^n(-1)^k=0\)，并将相邻两项配对得
\(\sum_{k=1}^n(-1)^k(k-1)=\frac n2\). 若 \(n\) 为奇数，则
\(\sum_{k=1}^n(-1)^k=-1\)，并得
\(\sum_{k=1}^n(-1)^k(k-1)=-\frac{n-1}{2}\). 因此
\(S_n=\begin{cases}
3^n-1+\dfrac n2\ln3,&n\text{为偶数},\\
3^n-1-\ln2-\dfrac{n-1}{2}\ln3,&n\text{为奇数}.
\end{cases}\)
\end{enumerate}
\end{solution}

\begin{problem}
某企业拟建造如图所示的容器（不计厚度，长度单位：米），其中容器的中间为圆柱形，左右两端均为半球形，按照设计要求容器的容积为 \(\frac{80\pi}{3}\text{ 立方米}\)，且 \(l\geqslant2r\). 假设该容器的建造费用仅与其表面积有关. 已知圆柱形部分每平方米建造费用为 \(3\text{ 千元}\)，半球形部分每平方米建造费用为 \(c\text{ 千元}\)（\(c>3\)）. 设该容器的建造费用为 \(y\text{ 千元}\).
\begin{enumerate}
\item 写出 \(y\) 关于 \(r\) 的函数表达式，并求该函数的定义域；
\item 求该容器的建造费用最小时的 \(r\).
\end{enumerate}

\bitmapfigure[max width=0.42\linewidth]{img/2011/shandong_science/q21_fig1.png}
\end{problem}

\begin{solution}
\begin{enumerate}
\item 圆柱部分的容积为 \(\pi r^2l\)，两个半球合起来是一个球，容积为 \(\frac43\pi r^3\). 因此
\(\pi r^2l+\frac43\pi r^3=\frac{80\pi}{3}\)，解得
\(l=\frac{80}{3r^2}-\frac43r\).
由设计条件 \(l\geqslant2r\)，得
\(\frac{80}{3r^2}-\frac43r\geqslant2r\)，即 \(80\geqslant10r^3\). 又 \(r>0\)，所以 \(0<r\leqslant2\).
圆柱侧面积为 \(2\pi rl\)，两端半球的总表面积为一个球的表面积 \(4\pi r^2\). 因而费用函数为
\(y=3\cdot2\pi rl+c\cdot4\pi r^2\)
\(=6\pi r\left(\frac{80}{3r^2}-\frac43r\right)+4c\pi r^2
=\frac{160\pi}{r}+4\pi(c-2)r^2\).
所以 \(y\) 的定义域为 \(0<r\leqslant2\).
\item 在 \(0<r\leqslant2\) 上，
\(y'(r)=8\pi(c-2)r-\frac{160\pi}{r^2}
=\frac{8\pi}{r^2}\bigl((c-2)r^3-20\bigr)\).
若 \(3<c\leqslant\frac92\)，则对 \(0<r\leqslant2\) 有
\((c-2)r^3-20\leqslant8(c-2)-20\leqslant0\)，所以 \(y\) 在定义域上递减，最小值在 \(r=2\) 处取得.
若 \(c>\frac92\)，令 \(y'(r)=0\)，得
\(r_0=\sqrt[3]{\frac{20}{c-2}}<2\). 当 \(0<r<r_0\) 时 \(y'(r)<0\)，当 \(r_0<r\leqslant2\) 时 \(y'(r)>0\)，所以最小值在 \(r=r_0\) 处取得.
综上，\(r=2\)（\(3<c\leqslant\frac92\)），或
\(r=\sqrt[3]{\frac{20}{c-2}}\)（\(c>\frac92\)）.
\end{enumerate}
\end{solution}

\begin{problem}
已知动直线 \(l\) 与椭圆 \(C: \frac{x^2}{3} + \frac{y^2}{2} = 1\) 交于 \(P(x_1, y_1)\)，\(Q(x_2, y_2)\) 两不同点，且 \(\triangle OPQ\) 的面积 \(S_{\triangle OPQ} = \frac{\sqrt{6}}{2}\). 其中 \(O\) 为坐标原点.
\begin{enumerate}
\item 证明： \(x_1^2 + x_2^2\) 和 \(y_1^2 + y_2^2\) 均为定值；
\item 设线段 \(PQ\) 的中点为 \(M\)，求 \(|OM| \cdot |PQ|\) 的最大值；
\item 椭圆 \(C\) 上是否存在三点 \(D\)，\(E\)，\(G\)，使得 \(S_{\triangle DOE} = S_{\triangle ODG} = S_{\triangle OEG} = \frac{\sqrt{6}}{2}\)；若存在，判断 \(\triangle DEG\) 的形状；若不存在，请说明理由.
\end{enumerate}
\end{problem}

\begin{solution}
\begin{enumerate}
\item 作伸缩变换
\(X=\dfrac{x}{\sqrt3}\)、\(Y=\dfrac{y}{\sqrt2}\). 椭圆 \(C\) 上的点分别变为单位圆上的点
\(P'=(X_1,Y_1)\)、\(Q'=(X_2,Y_2)\)，因为
\(X_i^2+Y_i^2=1\).
由三角形面积公式，
\(\dfrac12|x_1y_2-x_2y_1|=\dfrac{\sqrt6}{2}\)，所以
\(|x_1y_2-x_2y_1|=\sqrt6\). 又
\(x_1y_2-x_2y_1=\sqrt6(X_1Y_2-X_2Y_1)\)，故
\(|X_1Y_2-X_2Y_1|=1\).
对于两个单位向量，有恒等式
\((X_1Y_2-X_2Y_1)^2+(X_1X_2+Y_1Y_2)^2
=(X_1^2+Y_1^2)(X_2^2+Y_2^2)=1\).
因此 \(X_1X_2+Y_1Y_2=0\)，即 \(P'\)、\(Q'\) 互相垂直. 两个互相垂直的单位向量组成平面内的一组标准正交基，所以
\(X_1^2+X_2^2=1\)、\(Y_1^2+Y_2^2=1\). 于是
\(x_1^2+x_2^2=3(X_1^2+X_2^2)=3\)，
\(y_1^2+y_2^2=2(Y_1^2+Y_2^2)=2\)，均为定值.
\item 由第一问，\(P'\)、\(Q'\) 是单位圆上相互垂直的两点. 先取
\(P'=(\cos\theta,\sin\theta)\)、\(Q'=(-\sin\theta,\cos\theta)\). 另一种方向只会使下式中的 \(\sin2\theta\) 变号，结论相同. 记 \(t=\sin2\theta\)，则对应的椭圆上的点为
\(P=(\sqrt3\cos\theta,\sqrt2\sin\theta)\)，
\(Q=(-\sqrt3\sin\theta,\sqrt2\cos\theta)\).
中点 \(M\) 满足
\(|OM|^2=\dfrac14\left[3(\cos\theta-\sin\theta)^2+2(\sin\theta+\cos\theta)^2\right]=\dfrac{5-t}{4}\)，
而
\(|PQ|^2=3(\cos\theta+\sin\theta)^2+2(\sin\theta-\cos\theta)^2=5+t\).
所以
\((|OM|\cdot|PQ|)^2=\dfrac{(5-t)(5+t)}4=\dfrac{25-t^2}{4}\leqslant\dfrac{25}{4}\).
当 \(t=0\) 时等号成立，例如取 \(\theta=0\)，此时两点仍满足题设面积条件. 故
\(\max |OM|\cdot|PQ|=\dfrac52\).
\item 假设椭圆上存在满足条件的三点 \(D\)，\(E\)，\(G\). 对它们作同一伸缩变换，得到单位圆上的非零向量 \(D'\)、\(E'\)、\(G'\). 例如
\(S_{\triangle ODE}=\dfrac12|\det(D,E)|=\dfrac{\sqrt6}{2}\)，而伸缩变换使行列式乘以 \(\sqrt6\)，所以
\(|\det(D',E')|=1\). 由于 \(D'\)、\(E'\) 都是单位向量，这说明它们互相垂直. 同理，另外两个面积条件分别说明
\(D'\perp G'\)、\(E'\perp G'\).
但平面内不可能存在三个两两垂直的非零向量，而单位圆上的三个向量均非零，矛盾. 因此不存在这样的三点 \(D\)，\(E\)，\(G\)，也就没有需要判断形状的三角形.
\end{enumerate}
\end{solution}
